\section{Computational study on synthetic instances}\label{sec:experiments}

This section measures the algorithms and checks the results of Sections~\ref{sec:structure}
to~\ref{sec:knowledge} on instances drawn at random with fixed seeds. \emph{None of its numbers describes the repair
centre.} Monetary amounts are in arbitrary units. Every table and figure is regenerated by one command
(Appendix~\ref{app:repro}); computing times are those of the machine on which the command was last run.

\subsection{Exactness and cost of pricing}

The implementation solves \eqref{eq:P} by bisection on the labour price, pricing each type by
Remark~\ref{rem:twolevel}. In the test suite its optimal value equals that of an exact linear-programming solver
within $10^{-6}$ on 600 random instances, including instances with missing data, a negative treatment cost and a
rough opening longer than the careful one, and on unit types of nine and ten components where the pricing is done
without enumeration. Table~\ref{tab:pricing} and Figure~\ref{fig:pricing} compare the pricing of one unit type, at
\SynPrixValeurs{} values of the labour price, by enumeration of its plans and by Proposition~\ref{prop:pricing}.
The two agree to $\SynPrixEcart$ on every instance where both run. Enumeration grows by a factor of three per
component and is out of reach beyond a dozen; the separable pricing handles a unit of \SynPrixGrand{} components in
\SynPrixGrandTemps{} milliseconds, most of which is set-up.

\begin{table}[H]
\centering\small
\caption{Pricing one unit type at \SynPrixValeurs{} labour prices: time per type, set-up included (milliseconds).
Synthetic instances.}\label{tab:pricing}
\begin{tabular}{rrrr}
\toprule
Components & Plans & Enumeration & Proposition~\ref{prop:pricing} \\
\midrule
\corps{t_prix}
\bottomrule
\end{tabular}
\end{table}

\begin{figure}[H]
\centering
\includegraphics[width=0.48\textwidth]{\genere/fig_prix.pdf}\hfill
\includegraphics[width=0.48\textwidth]{\genere/fig_variabilite.pdf}
\caption{Left: time to price one unit type, by enumeration and by Proposition~\ref{prop:pricing} (logarithmic
scales; values in Table~\ref{tab:pricing}). Right: variability gap $\Gamma$ as a share of the fluid value, against
the dispersion of weekly arrivals, for three levels of available hours (values in
Table~\ref{tab:gap}).}\label{fig:pricing}
\end{figure}

\subsection{How much can a dynamic policy gain?}

Proposition~\ref{prop:fluid} bounds the gain of any dynamic policy by the variability gap $\Gamma$.
Table~\ref{tab:gap} computes its sample version on a line of \SynGammaTypes{} unit types over
\SynGammaSemaines{} weeks, for arrivals of the same mean and increasing dispersion, and for hours equal to 0.5, 0.8
and 1.2 times the hours the average flow would use without a limit. With a constant flow the gap is zero, as it
must be. It grows with dispersion and falls as hours become plentiful, and it reaches \SynGammaMax{} of the fluid
value in the most dispersed, most constrained case. The reading for a manager is direct: with Poisson-like
arrivals and hours close to the need, a buffer or a forecast has little to recover; with bursty arrivals and
scarce hours, it has much, and the gap says how much before any dynamic policy is built.

\begin{table}[H]
\centering\small
\caption{Variability gap $\Gamma$ as a share of the fluid value $V(\bar A,H)$, by arrival law and available hours
(as a multiple of the unconstrained need at the average flow). Synthetic instance.}\label{tab:gap}
\begin{tabular}{lrrrr}
\toprule
 & Coefficient of variation & \multicolumn{3}{c}{Hours / need} \\
\cmidrule(lr){3-5}
Weekly arrivals & of total arrivals & 0.5 & 0.8 & 1.2 \\
\midrule
\corps{t_variabilite}
\bottomrule
\end{tabular}
\end{table}

\subsection{Learning a yield that the prior gets wrong}

Table~\ref{tab:learning} and Figure~\ref{fig:regret} report a one-component benchmark in which the prior mean of the
yield is \SynTSprior, the true yield \SynTSvrai{} and the reuse threshold of Proposition~\ref{prop:threshold}
\SynTSseuil: reuse is the right decision and the prior says otherwise. The rule that follows the prior and the
greedy rule that follows the posterior mean never test the component, never learn, and lose the same amount on
every copy. Posterior sampling tests almost every copy, recovers the yield and keeps a regret that flattens. The
benchmark is deliberately minimal: it isolates the failure of greedy rules that Section~\ref{sec:learning}
describes. A benchmark on full instances of \eqref{eq:P} is left for future work.

\begin{table}[H]
\centering\small
\caption{One-component learning benchmark (true yield \SynTSvrai, prior mean \SynTSprior, threshold
\SynTSseuil; the oracle earns \SynTSoracle{} per copy). Synthetic.}\label{tab:learning}
\begin{tabular}{lrrr}
\toprule
Policy & Regret per copy & Copies tested & Final estimate of the yield \\
\midrule
\corps{t_apprentissage}
\bottomrule
\end{tabular}
\end{table}

\begin{figure}[H]
\centering
\includegraphics[width=0.48\textwidth]{\genere/fig_regret.pdf}
\caption{Cumulative regret on the benchmark of Table~\ref{tab:learning}. The greedy rule coincides with the rule
that follows the prior: neither ever tests the component.}\label{fig:regret}
\end{figure}

\subsection{Fusion and the reliability of sources}

Table~\ref{tab:fusion} compares the simulated error of the fusion rule with the bound of
Proposition~\ref{prop:fusion}, for \SynFusionClasses{} possible classes. The bound is valid throughout and
conservative by up to an order of magnitude, the price of the union over classes; it is vacuous for one or three
assertions of a mediocre source, which is the regime in which the chain sends a component to review rather than
accept it. For an error bound of 2\%, a source of reliability 0.9 needs \SynFusionKhaut{} independent assertions and
a source of reliability 0.6 needs \SynFusionKbas.

\begin{table}[H]
\centering\small
\caption{Probability that the most probable class is not the true one after $K$ independent assertions of a source
of known reliability $r$ (\SynFusionClasses{} classes; ties count as errors; \SynFusionEssais{} simulated cases per
row), and the bound of Proposition~\ref{prop:fusion}.}\label{tab:fusion}
\begin{tabular}{rrrr}
\toprule
$r$ & $K$ & Simulated error & Bound \eqref{eq:rate} \\
\midrule
\corps{t_fusion}
\bottomrule
\end{tabular}
\end{table}

For Proposition~\ref{prop:strata}, consider a source with \SynStrateNacc{} pieces of evidence on accepted
components, of accuracy \SynStrateRacc, and \SynStrateNfile{} on queued components, of accuracy \SynStrateRfile:
its true accuracy is \SynStrateVraie. An audit validates \SynStrateKacc{} of the former at random, and the
reviewers \SynStrateKfile{} of the latter. Counting all validations together gives \SynStrateEnsemble{} on average
(the proposition predicts \SynStrateAttendu), an error of root mean square \SynStrateEqmEnsemble: the source would
be judged much worse than it is, and components it classifies correctly would be sent to review for nothing. The
stratified estimator gives \SynStrateStrates{} on average, with a root-mean-square error of \SynStrateEqmStrates.

\subsection{Where the acquisition effort goes}

Table~\ref{tab:chain} runs the acquisition chain on a synthetic bill of materials of \SynChainePieces{} components
whose annual masses are log-normal, so that a few components carry most of the stake, with simulated sources of
fixed reliability: a rule on the description, a simulated language model reading descriptions and documents, and a
retrieval service that finds a document for some components and a manufacturer's declaration for a few. Costs are
the working assumptions of the implementation: \CoutDescription{} per reading of a description, \CoutCollecte{}
per retrieval, \CoutExtraction{} per reading of a document, \CoutControle{} per verdict, and \CoutRevue{} per human
review; the acceptance threshold is \SeuilAcceptation.

The chain makes \SynChaineAppels{} model calls, against \SynChaineExhaustif{} for a pipeline that reads, retrieves
and checks every component (\SynChaineEconomie{} fewer), for a total cost of \SynChaineCout{} against a stake
ceiling of \SynChaineEnjeu. No component is charged more than its ceiling: the highest ratio of
spending to ceiling is \SynChaineRatio{} (Lemma~\ref{lem:budget}). It accepts \SynChaineAcceptees{} classes without review, of which
\SynChaineJustesse{} are right, and they carry \SynChaineEnjeuAcceptees{} of the stake. It sends
\SynChaineRevue{} components to a human, which carry \SynChaineEnjeuRevue{} of the stake, and leaves
\SynChaineLaissees{} whose class would earn less than a review costs. The table shows the mechanism: effort follows stake. In the top fifth, \SynChaineHautTraites{} of the
components are settled automatically or sent to a reviewer, and the rest are components whose class, as currently
believed, would earn less than a review costs. In the bottom fifth, \SynChaineBasLaissees{} are left, at an average
cost of \SynChaineBasCout{} per component, and nobody is asked. With these costs the scarce resource is not the
language model, whose calls cost cents, but the human reviewer, and the stake is what rations it.

\begin{table}[H]
\centering\small
\caption{The acquisition chain on a synthetic bill of materials, by fifth of components ranked by stake ceiling
(first fifth: highest). Shares of the fifth's components that are accepted without review, sent to human review,
or left without a class.}\label{tab:chain}
\begin{tabular}{crrrrrrr}
\toprule
 & Share of & Median & Model calls & Cost per & & Sent to & \\
Fifth & total stake & ceiling & per component & component & Accepted & review & Left \\
\midrule
\corps{t_chaine}
\bottomrule
\end{tabular}
\end{table}
